\documentclass[11pt]{article}
\usepackage[margin=1.1in]{geometry}
\usepackage{amsmath,amssymb,microtype,xcolor,mdframed}
\usepackage[colorlinks=true,linkcolor=blue!50!black,urlcolor=blue!50!black]{hyperref}
\newmdenv[backgroundcolor=blue!4,linecolor=blue!35,linewidth=0.8pt,
  skipabove=9pt,skipbelow=9pt,innerleftmargin=10pt,innerrightmargin=10pt]{keyidea}
\setlength{\emergencystretch}{2em}
\title{A Reciprocal Test for Collatz Correction:\\
\Large A Friendly Guide}
\author{Companion to ``Bounded Collatz Correction and Reciprocal Orbit Sums''\thanks{Written with Codex (OpenAI). The technical paper and Lean source contain the exact proofs. No claim of literature priority is made.}}
\date{September 5, 2026}
\begin{document}
\maketitle
\section{The result in ordinary language}
Start with a positive whole number. At an even number, divide by two. At an odd number, multiply by three, add one, and divide by two. Call the resulting sequence $n_0,n_1,n_2,\ldots$.

Each odd step includes a small addition. The repository tracks the accumulated effect of those additions after removing the multiplicative factors. Call this normalized correction $c_t$. The existing formal name \texttt{Supercritical} means that $c_t$ stays bounded.

There is now an exact test for that property. Add the reciprocals of all the numbers visited:
\[
\frac1{n_0}+\frac1{n_1}+\frac1{n_2}+\cdots.
\]
\begin{keyidea}
\textbf{Lean proves both directions.} The normalized correction stays bounded exactly when this infinite reciprocal sum is finite. This holds for every positive starting number, without assuming a favorable density of odd or even steps.
\end{keyidea}
Large orbit values contribute small amounts. Repeated visits to small values contribute a substantial amount each time. The sum therefore measures both how high an orbit travels and how much time it spends at lower heights.

For example, starting at one gives $1,2,1,2,\ldots$. Its reciprocal sum is $1+1/2+1+1/2+\cdots$, which diverges. The theorem therefore says its normalized correction is unbounded. There is no contradiction: boundedness of the orbit and boundedness of the normalized correction are different properties.

\newpage
\section{Why the proof works}
Write $S_L$ for the reciprocal sum through time $L-1$, and let $N$ be the starting number. Two finite inequalities do all the work:
\[
\begin{aligned}
N S_L&\le6c_L+N\left(\frac1{n_L}-\frac1N\right),\\
N+c_L&\le N e^{S_L}.
\end{aligned}
\]
The first says that a large reciprocal sum must be paid for by correction growth, apart from a bounded endpoint term. On an even step the reciprocal doubles, so its contribution is charged to the change in reciprocal size. On an odd step the correction increases enough to pay the charge. Summing cancels the intermediate endpoint terms.

The second says that a bounded reciprocal sum keeps correction under control. At each step the quantity $N+c_t$ grows by a factor no larger than $1+1/n_t$. The inequality $1+x\le e^x$ turns the product of these factors into an exponential of their sum.

\section{What this gives us, and what is still missing}
If the reciprocal sum converges, its individual terms tend to zero. The orbit must therefore eventually leave every finite range and never return. Lean proves this consequence as well.

The converse has not been proved: an orbit heading to infinity might grow slowly enough for the reciprocal sum to diverge. For arbitrary positive sequences this happens, for example, with $n_t=t+1$. That sequence is an illustration of a logical gap, not a proposed Collatz orbit.

\begin{keyidea}
\textbf{The next research question.} Can a positive integer Collatz orbit escape forever while its reciprocal sum diverges? The new equivalence identifies this remaining case precisely. It does not eliminate it, and it does not eliminate the bounded-correction case either.
\end{keyidea}
The implementation is \texttt{lean/Collatz/Reciprocal.lean}. Its main theorem is
\begin{center}\small\texttt{supercritical\_iff\_summable\_reciprocal}.\end{center}
The technical paper gives the proofs, exact assumptions, and rebuild commands. The expanded 48-declaration audit reports only standard foundational axioms for the selected results. Collatz remains unproved; this result makes one part of the proof program more precise.
\end{document}
